% Capítulo común a todos los manuales.
\chapter{Primeros pasos}\label{cap:primeros}

\section{¿Qué es DIGAIPE?}

DIGAIPE (Dirección de Gestión de Admisión, Ingreso y Permanencia Estudiantil) es el sistema
interno de la Universidad de Los Andes para manejar las admisiones de
estudiantes nuevos: desde que llega el listado de asignados hasta que cada estudiante queda
matriculado, con sus reportes y estadísticas.

Cada persona trabaja con un \textbf{rol}. El rol decide qué opciones del menú ve y qué puede hacer.
El trabajo pasa de un rol al siguiente en este orden:

\begin{center}
\begin{tikzpicture}[
  caja/.style={draw=ulaazul,line width=1pt,rounded corners=3pt,fill=white,
    font=\sffamily\footnotesize,align=center,text width=2.35cm,minimum height=1.45cm},
  flecha/.style={-{Stealth[length=6pt]},line width=1pt,color=ulaazul}]
  \node[caja] (a) {\textbf{Administrador}\\crea usuarios y\\asigna roles};
  \node[caja,right=0.45cm of a] (b) {\textbf{Carga de}\\\textbf{admisiones}\\proceso y listado};
  \node[caja,right=0.45cm of b] (c) {\textbf{Generación de}\\\textbf{planillas}\\imprime la planilla};
  \node[caja,right=0.45cm of c] (d) {\textbf{Matrícula y}\\\textbf{devolución}\\revisa documentos};
  \node[caja,right=0.45cm of d] (e) {\textbf{Estadísticas}\\reportes y\\cifras};
  \foreach \i/\j in {a/b,b/c,c/d,d/e} \draw[flecha] (\i) -- (\j);
  \node[caja,below=0.55cm of c,text width=7cm,minimum height=0.9cm,draw=ulagris,dashed]
    (f) {\textbf{Auditor}: puede ver todo lo anterior, sin modificar nada};
\end{tikzpicture}
\end{center}

\begin{tabularx}{\linewidth}{@{}>{\leavevmode\color{ulaazuloscuro}\sffamily\bfseries}p{4.2cm} X@{}}
  \toprule
  \normalfont\bfseries Rol & \bfseries Qué hace \\
  \midrule
  Administrador & Todo lo del sistema. Además crea y elimina usuarios, asigna roles, activa o
                  desactiva procesos, elimina lotes en borrador y consulta la auditoría. \\
  Carga de admisiones & Crea los procesos de admisión con sus fechas, carga el listado de
                  asignados (Excel) y mantiene la tabla de carreras. \\
  Generación de planillas & Busca a cada estudiante, corrige sus datos de contacto y genera
                  su planilla en PDF. \\
  Matrícula y devolución & Revisa los documentos de cada estudiante y lo marca como
                  matriculado o le devuelve la solicitud. \\
  Estadísticas & Saca reportes por fechas, carreras, facultades y modalidades, en pantalla,
                  PDF o CSV. \\
  Auditor (solo lectura) & Consulta todo, incluido el registro de auditoría, sin poder cambiar
                  nada. \\
  \bottomrule
\end{tabularx}

\begin{nota}
Una misma persona puede tener varios roles. En ese caso verá en el menú las opciones de todos
ellos.
\end{nota}

\section{Entrar al sistema}

Abra el navegador (Chrome, Edge o Firefox) y escriba la dirección de DIGAIPE que le dio el
administrador. Verá esta pantalla:

\begin{captura}{login}{2732}{1600}
  \anotad{1}{1623}{806}{2423}{882}
  \anotad{2}{1623}{1011}{2423}{1093}
  \anotad{3}{1623}{1156}{2423}{1254}
\end{captura}

\begin{leyenda}
  \item[1] \textbf{Correo electrónico} con el que el administrador creó su cuenta.
  \item[2] \textbf{Contraseña.} Se escribe con puntos para que nadie la vea.
  \item[3] Pulse \boton{Iniciar sesión} para entrar.
\end{leyenda}

\begin{importante}[Si no puede entrar]
Revise que el correo esté bien escrito y que no tenga activadas las mayúsculas. Si olvidó la
contraseña, comuníquese con el administrador de DIGAIPE: el sistema no tiene opción para
recuperarla por correo.
\end{importante}

\section{La pantalla de inicio}

Al entrar verá la página de inicio. Arriba está la barra azul con el menú; solo aparecen las
opciones que su rol le permite usar.

\begin{captura}{inicio}{2732}{1600}
  \anota{1}{43}{11}{594}{167}
  \anota{2}{864}{43}{2462}{151}
  \anota{3}{2570}{22}{2690}{81}
  \anota{4}{1274}{551}{2570}{1453}
\end{captura}

\begin{leyenda}
  \item[1] \textbf{Logo de la ULA.} Haga clic aquí en cualquier momento para volver al inicio.
  \item[2] \textbf{Menú principal.} Al pasar el ratón sobre \menu{Admisiones} o
        \menu{Parametros} se abre una lista con más opciones.
  \item[3] \textbf{Salir.} Cierra su sesión. Úselo siempre al terminar, sobre todo en
        computadoras compartidas.
  \item[4] \textbf{Accesos rápidos} a las tareas más comunes.
\end{leyenda}

\subsection{El menú Admisiones}

Pase el ratón sobre \menu{Admisiones} para ver sus opciones:

\begin{captura}{menu-admisiones}{2732}{1000}[0.8\linewidth]
  \anota{1}{1145}{178}{1544}{918}
\end{captura}

\begin{tabularx}{\linewidth}{@{}>{\leavevmode\color{ulaazul}\sffamily\bfseries}l X@{}}
  \toprule
  \normalfont\bfseries Opción & \bfseries Para qué sirve \\
  \midrule
  Gestión de Aspirantes & Ver el listado de admitidos y cargar nuevos listados. \\
  Procesos de Admisión & Crear procesos y definir los días en que están activos. \\
  Generación de Planilla & Generar la planilla PDF de cada estudiante. \\
  Revisión de Documentos & Matricular o devolver solicitudes. \\
  \bottomrule
\end{tabularx}

\section{Consejos generales}

\begin{itemize}
  \item Todo lo que usted hace queda registrado con su nombre, la fecha y la hora. No comparta
        su cuenta con nadie.
  \item Si el administrador le cambia los roles, el sistema cierra su sesión. Vuelva a entrar
        y verá el menú actualizado.
  \item Si algo no funciona como dice este manual, consulte el capítulo
        \emph{Problemas frecuentes} al final, o comuníquese con el administrador.
\end{itemize}
